
    %%%%%%%%%%%%%%%%%%%%%%%%%%%%%%%%%%%%%%%%%%%%%%%%%%%%%%%%%%%%%%%%%%%%%%%%%%%%%%%%
    %%% Classe do documento
    \documentclass[12pt, addpoints]{exam}
    \UseRawInputEncoding

    %%%%%%%%%%%%%%%%%%%%%%%%%%%%%%%%%%%%%%%%%%%%%%%%%%%%%%%%%%%%%%%%%%%%%%%%%%%%%%%%
    % preambulo.tex
    %
    % Arquivo de configuração de packages para uso o LaTeX, conforme minhas
    % preferências e modelos pessoais.
    %
    % Para maiores informações, visite:
    %    
    %
    % NÃO ALTERE SE NÃO SOUBER O QUE ESTÁ FAZENDO!
    %%%%%%%%%%%%%%%%%%%%%%%%%%%%%%%%%%%%%%%%%%%%%%%%%%%%%%%%%%%%%%%%%%%%%%%%%%%%%%%%


    %%%%%%%%%%%%%%%%%%%%%%%%%%%%%%%%%%%%%%%%%%%%%%%%%%%%%%%%%%%%%%%%%%%%%%%%%%%%%%%%
    %%% Estruturas de controle:
    \input{static/utils/controles}

    %%%%%%%%%%%%%%%%%%%%%%%%%%%%%%%%%%%%%%%%%%%%%%%%%%%%%%%%%%%%%%%%%%%%%%%%%%%%%%%%
    %%% Compilação condicional em PDF:
    \input{static/utils/pdf}

    %%%%%%%%%%%%%%%%%%%%%%%%%%%%%%%%%%%%%%%%%%%%%%%%%%%%%%%%%%%%%%%%%%%%%%%%%%%%%%%%
    %%% Configurações de layout da página:
    \input{static/utils/layout}

    %%%%%%%%%%%%%%%%%%%%%%%%%%%%%%%%%%%%%%%%%%%%%%%%%%%%%%%%%%%%%%%%%%%%%%%%%%%%%%%%
    %%% Configurações para autores:
    \input{static/utils/autores}

    %%%%%%%%%%%%%%%%%%%%%%%%%%%%%%%%%%%%%%%%%%%%%%%%%%%%%%%%%%%%%%%%%%%%%%%%%%%%%%%%
    %%% Cabeçalho e rodapé:
    %%% Atenção! É MUITO DIFÍCIL ajustar apropriadamente os running headers
    %%% e running footers da primeira vez. Você pode confiar no LaTeX e deixar que
    %%% ele faça o ajuste automaticamente ou pode quebrar a cabeça e
    %%% tentar melhorar algo que o LaTeX já faz muito bem, mesmo que não fique
    %%% exatamente do jeito que você quer. O que você prefere? Se quiser ajustar
    %%% manualmente, descomente a linha abaixo e faça as configurações no arquivo
    %%% "utils/headings.tex".
    %\input{static/utils/headings}

    %%%%%%%%%%%%%%%%%%%%%%%%%%%%%%%%%%%%%%%%%%%%%%%%%%%%%%%%%%%%%%%%%%%%%%%%%%%%%%%%
    %%% Configuração para as fontes do documento:
    %%% Se você quiser usar fontes próprias ou do sistema, precisará ajustar as
    %%% configurações no arquivo "utils/fontes.tex".
    %%% ATENÇÃO: por padrão eu utilizo XeLaTeX com fontes proprietárias projetadas
    %%% por Matthew Butterick, adquiridas em https://mbtype.com, que não podem ser
    %%% distribuídas devido à licença de utilização. Se você usar XeLaTeX, você
    %%% DEVERÁ OBRIGATORIAMENTE ajustar as fontes no arquivo "utils/fontes.tex".
    \input{static/utils/fontes}

    %%%%%%%%%%%%%%%%%%%%%%%%%%%%%%%%%%%%%%%%%%%%%%%%%%%%%%%%%%%%%%%%%%%%%%%%%%%%%%%%
    %%% Sumário:
    \input{static/utils/toc}

    %%%%%%%%%%%%%%%%%%%%%%%%%%%%%%%%%%%%%%%%%%%%%%%%%%%%%%%%%%%%%%%%%%%%%%%%%%%%%%%%
    %%% Matemática:
    \input{static/utils/matematica}

    %%%%%%%%%%%%%%%%%%%%%%%%%%%%%%%%%%%%%%%%%%%%%%%%%%%%%%%%%%%%%%%%%%%%%%%%%%%%%%%%
    %%% Notas de rodapé:
    \input{static/utils/footnotes}

    %%%%%%%%%%%%%%%%%%%%%%%%%%%%%%%%%%%%%%%%%%%%%%%%%%%%%%%%%%%%%%%%%%%%%%%%%%%%%%%%
    %%% Referências cruzadas, links, citações:
    \input{static/utils/crossrefs}

    %%%%%%%%%%%%%%%%%%%%%%%%%%%%%%%%%%%%%%%%%%%%%%%%%%%%%%%%%%%%%%%%%%%%%%%%%%%%%%%%
    %%% Configurações para os hiperlinks em PDF:
    \input{static/utils/pdfconfig}

    %%%%%%%%%%%%%%%%%%%%%%%%%%%%%%%%%%%%%%%%%%%%%%%%%%%%%%%%%%%%%%%%%%%%%%%%%%%%%%%%
    %%% Glossário e índice remissivo:
    \input{static/utils/glossindex}

    %%%%%%%%%%%%%%%%%%%%%%%%%%%%%%%%%%%%%%%%%%%%%%%%%%%%%%%%%%%%%%%%%%%%%%%%%%%%%%%%
    %%% Referências bibliográficas:
    \input{static/utils/bibliografia}

    %%%%%%%%%%%%%%%%%%%%%%%%%%%%%%%%%%%%%%%%%%%%%%%%%%%%%%%%%%%%%%%%%%%%%%%%%%%%%%%%
    %%% Cores:
    \input{static/utils/cores}

    %%%%%%%%%%%%%%%%%%%%%%%%%%%%%%%%%%%%%%%%%%%%%%%%%%%%%%%%%%%%%%%%%%%%%%%%%%%%%%%%
    %%% Computação:
    \input{static/utils/computacao}

    %%%%%%%%%%%%%%%%%%%%%%%%%%%%%%%%%%%%%%%%%%%%%%%%%%%%%%%%%%%%%%%%%%%%%%%%%%%%%%%%
    %%% Gráficos:
    \input{static/utils/graficos}

    %%%%%%%%%%%%%%%%%%%%%%%%%%%%%%%%%%%%%%%%%%%%%%%%%%%%%%%%%%%%%%%%%%%%%%%%%%%%%%%%
    %%% Ícones:
    \input{static/utils/icones}

    %%%%%%%%%%%%%%%%%%%%%%%%%%%%%%%%%%%%%%%%%%%%%%%%%%%%%%%%%%%%%%%%%%%%%%%%%%%%%%%%
    %%% Blocos:
    \input{static/utils/blocos}

    %%%%%%%%%%%%%%%%%%%%%%%%%%%%%%%%%%%%%%%%%%%%%%%%%%%%%%%%%%%%%%%%%%%%%%%%%%%%%%%%
    %%% Boxes coloridos:
    \input{static/utils/colorbox}

    %%%%%%%%%%%%%%%%%%%%%%%%%%%%%%%%%%%%%%%%%%%%%%%%%%%%%%%%%%%%%%%%%%%%%%%%%%%%%%%%
    %%% Tabelas:
    \input{static/utils/tabelas}

    %%%%%%%%%%%%%%%%%%%%%%%%%%%%%%%%%%%%%%%%%%%%%%%%%%%%%%%%%%%%%%%%%%%%%%%%%%%%%%%%
    %%% Ambientes de listas:
    \input{static/utils/listas}

    %%%%%%%%%%%%%%%%%%%%%%%%%%%%%%%%%%%%%%%%%%%%%%%%%%%%%%%%%%%%%%%%%%%%%%%%%%%%%%%%
    %%% Ambientes floats e similares:
    \input{static/utils/floats}

    %%%%%%%%%%%%%%%%%%%%%%%%%%%%%%%%%%%%%%%%%%%%%%%%%%%%%%%%%%%%%%%%%%%%%%%%%%%%%%%%
    %%% Ajustes para a classe EXAM:
    \input{static/utils/exam}

    %%%%%%%%%%%%%%%%%%%%%%%%%%%%%%%%%%%%%%%%%%%%%%%%%%%%%%%%%%%%%%%%%%%%%%%%%%%%%%%%
    %%% Ajustes para textos:
    \input{static/utils/texto}

    %%%%%%%%%%%%%%%%%%%%%%%%%%%%%%%%%%%%%%%%%%%%%%%%%%%%%%%%%%%%%%%%%%%%%%%%%%%%%%%%
    %%% Ambientes, aliases, comandos e customizações em geral para serem
    %%% utilizadas nos documentos. Defina aqui qualquer customização específica
    %%% para seus documentos!
    \input{static/utils/customizacoes}

    %%%%%%%%%%%%%%%%%%%%%%%%%%%%%%%%%%%%%%%%%%%%%%%%%%%%%%%%%%%%%%%%%%%%%%%%%%%%%%%%
    %%% Ajuste do layout, espaçamento de linhas e etc.:
    \geometry{a4paper, portrait, left=2.0cm, right=2.0cm, top=2.0cm, bottom=2.0cm}
    \singlespacing

    %%%%%%%%%%%%%%%%%%%%%%%%%%%%%%%%%%%%%%%%%%%%%%%%%%%%%%%%%%%%%%%%%%%%%%%%%%%%%%%%
    %%% Configurações de cabeçalho e rodapé (não use fancyhdr pois ocorre conflito):
    \pagestyle{headandfoot}
    \runningheadrule
    \firstpageheader{}{}{}
    \runningheader{Sem disciplina}{}{Avaliação}
    \firstpagefooter{}{}{Página \thepage\ de \numpages}
    \runningfooter{}{}{Página \thepage\ de \numpages}
    \runningfootrule

    %%%%%%%%%%%%%%%%%%%%%%%%%%%%%%%%%%%%%%%%%%%%%%%%%%%%%%%%%%%%%%%%%%%%%%%%%%%%%%%%
    %%% Configurações para as propriedades do PDF:
    \hypersetup{
    hidelinks,           % Comente para web, descomente para impressão
    colorlinks=true,      % True para web, False para impressão
    pdftitle=TesteProva: Avaliação,
    pdfauthor=Professor,
    pdfsubject=Sem disciplina,
    pdfkeywords=['Sem tag'],
    pdfinfo={
        CreationDate={}, % Ex.: D:AAAAMMDDHH24MISS
        ModDate={}       % Ex.: D:AAAAMMDDHH24MISS
    }
    }
    \input{static/utils/pdfconfig}

    %%%%%%%%%%%%%%%%%%%%%%%%%%%%%%%%%%%%%%%%%%%%%%%%%%%%%%%%%%%%%%%%%%%%%%%%%%%%%%%%
    %%% Começa o documento
    \begin{document}

    %%%%%%%%%%%%%%%%%%%%%%%%%%%%%%%%%%%%%%%%%%%%%%%%%%%%%%%%%%%%%%%%%%%%%%%%%%%%%%%%
    %%% Folha de instruções padronizadas da para a prova
    \pdfbookmark[1]{Instruções}{instrucoes}
    \begin{figure}[H]
        \begin{center}
            \includegraphics[scale=0.3]{{static/imgs/logo-.png}}
        \end{center}	
    \end{figure}

    \vspace{-1cm}

    \begin{table}[H]
        \centering
        \begin{tabular}{|llll|}
            \hline
            \multicolumn{2}{|l|}{\textbf{Disciplina:} Sem disciplina} & 
            \multicolumn{1}{l|}{\multirow{3}{*}{\textbf{\makecell{Nota:\\ \,\\ \,}}}} & 
            \multirow{3}{*}{\textbf{\makecell{Coordenador:\\ \,\\ \,}}} \\ 
            \cline{1-2}
            \multicolumn{2}{|l|}{\textbf{Professor:} Professor \hspace{4.72cm} } & 
            \multicolumn{1}{l|}{} & \\ 
            \cline{1-2}
            \multicolumn{2}{|l|}{\textbf{Aluno:}} & 
            \multicolumn{1}{l|}{} & \\ 
            \hline
            \multicolumn{1}{|l|}{\textbf{Turma:} \hspace{6.3cm} } & 
            \multicolumn{1}{l|}{\textbf{Semestre:}} & 
            \multicolumn{2}{l|}{\textbf{Valor:} 10 pontos} \\ 
            \hline
            \multicolumn{1}{|l|}{\textbf{Data:}} & 
            \multicolumn{3}{l|}{\textbf{Avaliação:}} \\ 
            \hline
        \end{tabular}
    \end{table}

    \begin{center}
        \fbox{\setlength\fboxsep{0.3cm} \fbox{\parbox{15.4cm}{
            \begin{center}
                \textbf{INSTRUÇÕES PARA A REALIZAÇÃO DESTA AVALIAÇÃO:}
            \end{center}
            \begin{itemize}[leftmargin=*]
                \item Esta é a primeira avaliação da disciplina Sem disciplina, 
                    e engloba o conteúdo referente Sem tag.
                \item Esta avaliação contém 50 questões objetivas, \textbf{todas obrigatórias}.
                \item \textbf{Leia atentamente cada questão!} As questões objetivas terão uma,
                    e apenas uma única, resposta a ser marcada.
                \item \textbf{Desligue o celular} antes de começar. A avaliação é
                    \textbf{individual} e \textbf{sem consulta}.
                \item As questões podem ser respondidas, na prova, com lápis ou caneta. Ao final
                    da prova \textbf{{VOCÊ DEVERÁ PREENCHER O GABARITO COM CANETA
                    DE COR PRETA}} \textbf{\underline{OU AZUL}} para que seja feita a correção
                    automatizada através da leitura do padrão de respostas marcado no
                    gabarito.
                \item \textbf{CUIDADO ao preencher o gabarito} pois eles são \textbf{INDIVIDUAIS
                    e NOMEADOS} (cada estudante tem um gabarito próprio específico, com um
                    código de identificação único). \textbf{Questões rasuradas são anuladas}
                    pelo software de leitura do gabarito.
                \item Ao final da prova, \textbf{DEVOLVA A PROVA E O GABARITO} para o professor.
                \item Siga todas as normas de \textbf{Integridade Acadêmica} da disciplina.
                    Alunos que forem flagrados com qualquer espécie de ``cola'' ou trocando
                    informações com outros alunos terão suas avaliações recolhidas, as notas
                    zeradas, e a situação será encaminhada para a coordenação para a aplicação
                    das penalidades previstas pela universidade.
                \item Com 90 minutos de avaliação você terá tempo de até 2,min por questão,
                    em média.
                \item Boa avaliação!
            \end{itemize}
            \vspace{2cm}
        }}}
    \end{center}
    \newpage

    %%%%%%%%%%%%%%%%%%%%%%%%%%%%%%%%%%%%%%%%%%%%%%%%%%%%%%%%%%%%%%%%%%%%%%%%%%%%%%%%
    %%% Inicia o bloco de questões:
    \begin{questions}

    \newpage
    \question

O contador da figura abaixo é:

\begin{figure}

    \centering

    \includegraphics[width=0.5\linewidth]{static/imgs/defaul.png}

    \label{fig:enter-label}

\end{figure}
\begin{choices}
\choice isócrono
\choice assíncrono
\choice auto-sincronizado
\choice síncrono
\choice anisócrono
\end{choices}

\question

A derivada da função \( f(x) = x^x \) é igual a:
\begin{choices}
\choice \( x^x (\ln(x) + 1) \)
\choice \( x^x \)
\choice \( x^x \ln(x) \)
\choice \( x^x (\ln(x) + x) \)
\choice \( x x^{x-1} \)
\end{choices}

\question

Considere as seguintes tabelas em uma base de dados relacional (chaves primárias sublinhadas):\\

Departamento (\underline{CodDepto}, NomeDepto)\\

Empregado (\underline{CodEmp}, NomeEmp, CodDepto)\\

Considere as seguintes restrições de integridade sobre esta base de dados relacional:\\

- Empregado.CodDepto é sempre diferente de NULL\\

- Empregado.CodDepto é chave estrangeira da tabela Departamento com cláusulas ON DELETE RESTRICT e ON UPDATE RESTRICT\\

Qual das seguintes validações não é especificada por estas restrições de integridade:
\begin{choices}
\choice Sempre que uma linha for excluída de Departamento, deve ser garantido que o valor de Departamento.CodDepto não aparece na coluna Empregado.CodDepto.
\choice Sempre que o valor de Empregado.CodDepto for alterado, deve ser garantido que o novo valor de Empregado.CodDepto aparece em Departamento.CodDepto.
\choice Sempre que o valor de Departamento.CodDepto for alterado, deve ser garantido que não há uma linha com o antigo valor de Departamento.CodDepto na coluna Empregado.CodDepto.
\choice Sempre que uma nova linha for inserida em Departamento, deve ser garantido que o valor de Departamento.CodDepto aparece na coluna Empregado.CodDepto.
\choice Sempre que uma nova linha for inserida em Empregado, deve ser garantido que o valor de Empregado.CodDepto aparece na coluna Departamento.CodDepto.
\end{choices}

\question

Considere a seguinte tabela para uma base de dados relacional:\\

Empregado (CodEmp, NomeEmp, CodDepto)\\

Considere que esta tabela tem um índice na forma de uma árvore B sobre as colunas (CodEmp, CodDepto), nesta ordem.

Quanto a este índice, considere as seguintes afirmativas:

\begin{enumerate}

    \item Este índice pode ser usado pelo SGBD relacional para acelerar uma consulta na qual são fornecidos os valores de CodEmp e CodDepto.

    \item Este índice pode ser usado pelo SGBD relacional para acelerar uma consulta na qual é fornecido um valor de CodEmp.

    \item Este índice não é adequado para ser usado pelo SGBD relacional para acelerar uma consulta na qual é fornecido um valor de CodDepto.

    \item O algoritmo que faz inserções e remoções de entradas do índice tem por objetivo garantir que o índice fique organizado de tal forma que o acesso a cada nodo da árvore implique em número de acessos semelhantes.

    \item O índice por árvore-B não é adequado para tabelas que sofrem grande número de inclusões e exclusões, pois exige reorganizações frequentes.

\end{enumerate}





Quanto a estas afirmativas pode se dizer que:
\begin{choices}
\choice Apenas as afirmativas 1), 2), 3) e 4) estão corretas
\choice Todas afirmativas estão corretas
\choice Apenas as afirmativas 1), 2) e 5) estão corretas
\choice Apenas as afirmativas 1), 2) e 4) estão corretas
\choice Nenhuma das afirmativas está correta
\end{choices}

\question

Um agente SNMP é um aplicativo que é executado:
\begin{choices}
\choice em "firewalls" com o objetivo de proteger acesso a rede
\choice a partir de um computador específico para monitorar a rede
\choice em um dispositivo de rede
\choice em roteadores com filtragem de pacotes com o objetivo de proteger acesso a rede
\choice em computadores denominados de gerentes
\end{choices}

\question

Um sistema centralizado é um concentrador de recursos; um sistema distribuído apresenta seus recursos dispersos. Entretanto, nem todo conjunto de recursos computacionais dispersos pode ser considerado um sistema distribuído. Considerando um conjunto de computadores, assinale a alternativa que melhor corresponde às características necessárias para considerá-lo um sistema distribuído:
\begin{choices}
\choice Existência de memória secundária compartilhada e protocolos de sincronização de estado
\choice Existência de sistema operacional idêntico e hardware padronizado em todos os computadores
\choice Existência de memória compartilhada e relógios locais sincronizados 
\choice Suporte de rede e um relógio global
\choice Suporte de rede e funções primitivas de comunicação 
\end{choices}

\question

Uma árvore binária é declarada em C como



\begin{verbatim}

typedef struct no *apontador;

struct no {

    int valor;

    apontador esq, dir;

};

\end{verbatim}



onde \texttt{esq} e \texttt{dir} representam ligações para os filhos esquerdo e direito de um nó da árvore, respectivamente. Qual das seguintes alternativas é uma implementação correta da operação que inverte as posições dos filhos esquerdo e direito de um nó \texttt{p} da árvore, onde \texttt{t} é um apontador auxiliar.


\begin{choices}
\choice \texttt{t = p;} \\

          \texttt{p->esq = p->dir;} \\

          \texttt{p->dir = p->esq;}
\choice \texttt{t = p->dir;} \\

          \texttt{p->esq = p->dir;} \\

          \texttt{p->dir = t;}
\choice \texttt{t = p->dir;} \\

          \texttt{p->dir = p->esq;} \\

          \texttt{p->esq = t;}
\choice \texttt{p->esq = p->dir;} \\

          \texttt{t = p->esq;} \\

          \texttt{p->dir = t;}
\choice \texttt{p->dir = t;} \\

          \texttt{p->esq = p->dir;} \\

          \texttt{p->dir = t;}
\end{choices}

\question

Considere as seguintes afirmações sobre mecanismos de inferência em sistemas baseados em regras.

\begin{enumerate}[label=\Roman*.]

    \item O encadeamento regressivo tem pouca utilidade prática, pois deve partir do possível resultado.

    \item O encadeamento progressivo tanto pode ser em amplitude quanto em profundidade.

    \item Podem trabalhar com informações incertas ou incompletas. 

\end{enumerate}

São corretas:
\begin{choices}
\choice Apenas I e II
\choice Apenas I e III
\choice Apenas III
\choice I, II e III
\choice Apenas II e III
\end{choices}

\question

Considere as afirmações abaixo sobre endereços IP:

\begin{itemize}

    \item[I.] Uma rede IP classe C fornece até 256 endereços válidos para serem atribuídos a equipe.

    \item[II.] A quantidade máxima de bits que pode ser utilizada para se definir sub-redes em uma rede IP classe C é seis (6).

    \item[III.] A máscara padrão para uma rede classe B é 255.255.255.0.

\end{itemize}

Qual das seguintes alternativas representa as afirmações corretas?
\begin{choices}
\choice Somente I e II.
\choice Somente I.
\choice Somente II e III.
\choice Somente II.
\choice Somente III.
\end{choices}

\question

Uma característica de uma arquitetura RISC é:
\begin{choices}
\choice Possui um pequeno conjunto de instruções simples
\choice A arquitetura é constituída de milhares de processadores
\choice O processador é formado por válvulas e transistores
\choice Uma arquitetura de alto risco pois o mercado de hardware evolui muito rapidamente
\choice Possui um grande conjunto de instruções complexas
\end{choices}

\question

Dada a seguinte fórmula (lógica de primeira ordem):\\

$\forall x \ \exists y \ | \ ama(x,y)$ \\

qual das seguintes sentenças em linguagem natural ela representa, considerando que \( ama(x,y) \) representa que \( x \) ama \( y \)?


\begin{choices}
\choice Há alguém que todos amam.
\choice Ninguém ama a todos.
\choice Todos amam alguém.
\choice Alguém ama a todos.
\choice Nenhuma das anteriores.
\end{choices}

\question

Qual é a opção que descreve a tarefa executada pelo seguinte algoritmo escrito em Pascal? 

\begin{verbatim} 

procedure fazalgo (var x, var y)

begin 

	  x := x + y; 

	  y := x - y;

	  x := x - y; 

end 

\end{verbatim}
\begin{choices}
\choice divide y por x utilizando a subtração e retorna o resultado em x
\choice não altera os valores de x e y
\choice calcula o mínimo múltiplo comum entre x e y e retorna o valor em x
\choice troca os valores de x e y
\choice divide x por y utilizando a subtração e retorna o resultado em x
\end{choices}

\question

O usuário que não é da área de tecnologia, costuma ser de alta gerência, que é

especialista no negócio e tem a responsabilidade central pelas políticas de

dados da empresa é o:
\begin{choices}
\choice Coordenador de Tecnologia da Informação (ITS: IT Supervisor)
\choice Administrador de Dados (DA: Data Administrator)
\choice Administrador do Banco de Dados (DBA: Database Administrator)
\choice Encarregado da Proteção de Dados (DPO: Data Protection Officer)
\choice Gerente de Tecnologia da Informação (ITM: IT Manager)
\end{choices}

\question

Considere um problema em que são dados 5 objetos com os seguintes pesos e valores:\\

\text{pesos: } $(W_1, W_2, W_3, W_4, W_5) = (6, 10, 9, 5, 12)$\\

\text{valores: } $(P_1, P_2, P_3, P_4, P_5) = (8, 5, 10, 15, 7)$\\

Além disso, é dada uma mochila que suporta até 30 unidades de peso, para transportar os objetos. O objetivo do problema é preencher a mochila de tal forma que o valor total dos objetos a serem transportados seja o maior possível, mas sem exceder o limite de peso suportado pela mochila. Assuma que é permitido colocar fração de um objeto na mochila. Qual das seguintes alternativas corresponde ao valor máximo obtido no preenchimento da mochila:
\begin{choices}
\choice 30.34
\choice 21.5
\choice 43.1
\choice 38.83
\choice 12.2
\end{choices}

\question

Qual das afirmações a seguir, relativas à análise sintática, está INCORRETA?


\begin{choices}
\choice Algoritmos descendentes recursivos podem ser utilizados para algumas gramáticas ambíguas.
\choice Algoritmos de análise redutiva podem ser utilizados mesmo para gramáticas ambíguas.
\choice Analisadores sintáticos descendentes recursivos são mais simples de implementar do que analisadores sintáticos redutivos.
\choice Uma das diferenças entre os diversos algoritmos de análise redutiva é a forma de identificar o \textit{handle} na pilha.
\choice As gramáticas LL podem descrever mais linguagens do que as gramáticas LR.
\end{choices}

\question

Professor Mac Sperto propôs o seguinte algoritmo de ordenação, chamado de Super Merge, similar ao \textit{merge sort}: divida o vetor em 4 partes do mesmo tamanho (ao invés de 2, como é feito no \textit{merge sort}). Ordene recursivamente cada uma das partes e depois intercale-as por um procedimento semelhante ao procedimento de intercalação do \textit{merge sort}. Qual das alternativas abaixo é verdadeira?
\begin{choices}
\choice Super Merge está correto, mas consome tempo \( O(\text{merge sort}) \)
\choice Nenhuma das afirmações acima está correta
\choice Super Merge não está correto. Não é possível ordenar quebrando o vetor em 4 partes
\choice Super Merge está correto, mas consome tempo maior que \( O(\text{merge sort}) \)
\choice Super Merge está correto, mas consome tempo menor que \( O(\text{merge sort}) \)
\end{choices}

\question

Se \( \lim_{n \to \infty} \left( \frac{1}{n^2} + \frac{2}{n^2} + \cdots + \frac{n-1}{n^2} \right) = L \) então


\begin{choices}
\choice \( L = \infty \)
\choice \( L = 1/2 \)
\choice \( L = 0 \)
\choice \( L = 1 \)
\choice \( L = 2 \)
\end{choices}

\question

Em relação aos metadados de um banco de dados, assinale a afirmação correta:
\begin{choices}
\choice Armazenam informações sobre os dados dos usuários, ou seja, são dados de rotina dos usuários, aqueles dados que os usuários precisam para trabalhar no seu dia a dia.
\choice Armazenam informações como os dados de cadastro dos usuários, os valores de compra em uma ordem de compra, as turmas e disciplinas que os professores dão aula, e dados de recursos humanos.
\choice São atualizados pelo SGBD, sempre que ocorre alguma operação de inserção, atualização ou remoção de dados.
\choice São armazenados no catálogo (ou dicionário de dados) do sistema de gerenciamento de bancos de dados, que fica nas mesmas tabelas de dados dos usuários, permitindo assim que os metadados estejam sempre atualizados.
\choice Como são sempre atualizados de forma automática pelo sistema de gerenciamento de bancos de dados, estão sempre prontos para nos dar uma ``foto' da estrutura do banco de dados em um momento no tempo.
\end{choices}

\question

Qual é a função do circuito da figura abaixo?

\begin{figure}

    \centering

    \includegraphics[width=0.5\linewidth]{static/imgs/defaul.png}

    \label{fig:enter-label}

\end{figure}
\begin{choices}
\choice somador
\choice deslocador
\choice subtrator
\choice multiplexador
\end{choices}

\question

Considerando a rede de Petri abaixo, quais das alternativas são verdadeiras?



\begin{itemize}

    \item[I)] O lugar \( A \) está habilitado a disparar.

    \item[II)] Apenas a transição \( T1 \) está habilitada a disparar.

    \item[III)] A sequência de transições \( (T1, T2, T3, T2) \) pode ser disparada, nessa ordem.

    \item[IV)] A transição \( T4 \) nunca poderá ser disparada.

\end{itemize}



\begin{figure}[H]

    \centering

    \includegraphics[width=0.5\linewidth]{static/imgs/defaul.png}

    \caption{Questão 49}

    \label{fig:enter-label}

\end{figure}
\begin{choices}
\choice Apenas as alternativas II e III.
\choice Apenas as alternativas I e III.
\choice Apenas as alternativas II, III e IV.
\choice Todas as alternativas.
\choice Apenas as alternativas III, IV.
\end{choices}

\question

Na definição da aquisição de um novo software de

gerenciamento de bancos de dados (SGBD) para uma empresa da área de transporte 

coletivo urbano, a direção da área de informática conduziu o processo de decisão

da seguinte forma: foi designado um profissional da área de bancos de dados

(aquele com maior experiência na área) e atribuída a ele a tarefa de decidir

qual seria o melhor SBGD a ser adquirido. Esse profissional desenvolveu uma

série de estudos sobre as opções disponíveis, utilizando ténicas de simulação e

testes específicos, levando em conta que a quase totalidade dos dados que serão

armazenados são dados de cadastro altamente estruturados e que devem estar

disponíveis na Internet para consulta pelos usuários. Ao final, apresentou

ao diretor um relatório em que indicava claramente qual o melhor tipo de SGBD a

ser adquirido pela empresa. Com base nessa informação, o diretor da empresa

disparou o processo de compra do SGBD. Podemos concluir que o profissional de

bancos de dados escolheu um:
\begin{choices}
\choice SGBD Grafo, para armazenar todos os relacionamentos entre os usuários do transporte público nos ônibus.
\choice SGBD Não Relacional (NoSQL) já que, para armazenar todos os dados e ainda permitir a consulta pelos usuários na Internet, o melhor seria utilizar uma solução com bancos de dados não relacionais que são otimizados para uso na web.
\choice Nenhuma das respostas acima.
\choice SGBD Relacional, para armazenar de forma otimizada os dados que, em sua maioria, são altamente estruturados.
\choice SGBD Hierárquico, para armazenar corretamente e de forma otimizada todas os relacionamentos entre os ônibus, empregados e passageiros do transporte público.
\end{choices}

\question

Considere o seguinte trecho de programa:\\

1. \( i := 1; \)\\

2. while \( i \leq n \) do\\

   begin\\

3. \( sum := sum + a[i]; \)\\

4. \( i := i + 1; \)\\

   end;\\



Considere que:\\

- \( I \) representa a inicialização da variável \( i := 1 \) na linha 1;\\

- \( T \) representa o teste da linha 2;\\

- \( A \) representa os comandos da linha 3;\\

- \( P \) representa o incremento na linha 4.\\



Qual é a expressão regular que representa todas as sequências de passos possíveis de serem executados por este trecho de programa?
\begin{choices}
\choice \( I(TAP)^* \)
\choice \( I(T(AP)^*)^+ \)
\choice \( I(T(AP)^*)^* \)
\choice \( I(TAP)^+ \)
\choice \( IT^+A^*P^* \)
\end{choices}

\question

Dado um vetor \( u \in \mathbb{R}^2 \), \( u = (-3, 4) \), vamos denotar por \( v \) o vetor de \( \mathbb{R}^2 \) que tem tamanho 1 e é ortogonal a \( u \). Então \( v \) pode ser dado por


\begin{choices}
\choice \( \left( \frac{3}{5}, \frac{4}{5} \right) \)
\choice \( \left( -\frac{4}{5}, \frac{2}{5} \right) \)
\choice \( \left( -\frac{4}{5}, -\frac{3}{5} \right) \)
\choice \( \left( -\frac{4}{5}, \frac{3}{5} \right) \)
\choice \( \left( -\frac{4}{5}, \frac{1}{5} \right) \)
\end{choices}

\question

Em qual das situações abaixo um sistema de Raciocínio Baseado em Casos não deve ser utilizado?


\begin{choices}
\choice Em aplicações de diagnóstico médico.
\choice Quando especialistas conversam sobre seus domínios dando exemplos.
\choice Quando for fácil a obtenção de regras
\choice Quando a experiência for tão valiosa quanto o conhecimento em livros texto.
\choice Quando as regras utilizadas apresentam um grande número de exceções.
\end{choices}

\question

Considere a seguinte tabela em uma base de dados relacional (chave primária sublinhada):\\

Tabela1 (CodAluno, CodDisciplina, AnoSemestre, NomeAluno, NomeDisciplina, CodNota, DescricaoNota)\\

Considere as seguintes dependências funcionais:

\begin{itemize}

    \item CodAluno \(\rightarrow\) NomeAluno

    \item CodDisciplina \(\rightarrow\) NomeDisciplina

    \item (CodAluno, CodDisciplina, AnoSemestre) \(\rightarrow\) CodNota

    \item (CodAluno, CodDisciplina, AnoSemestre) \(\rightarrow\) DescricaoNota

    \item CodNota \(\rightarrow\) DescricaoNota

\end{itemize}



Considerando as formas normais, qual das afirmativas abaixo se aplica:
\begin{choices}
\choice A tabela encontra-se na segunda forma normal, mas não na terceira forma normal.
\choice A tabela está na quarta forma normal.
\choice A tabela encontra-se na terceira forma normal, mas não na quarta forma normal.
\choice A tabela encontra-se na primeira forma normal, mas não na segunda forma normal.
\choice A tabela não está na primeira forma normal.
\end{choices}

\question

(ENADE 2014) Considere o seguinte banco de dados ilustrado no diagrama relaconal

abaixo:

\begin{figure}[H]

  \centering

  \caption{Estados e fornecedores}

  \label{fig:estados}

  %\fbox{

    \includegraphics[width=0.5\linewidth]{static/imgs/defaul.png}

  %}\\

  %\footnotesize{Fonte: xxxx.}

\end{figure}

A expressão SQL que obtém os nomes dos estados para os quais não há fornecedores

cadastrados é:


\begin{choices}
\choice \begin{verbatim}SELECT e.nome

FROM estado e,

FROM fornecedor f

WHERE e.nome = f.uf;

\end{verbatim}


\choice \begin{verbatim}SELECT e.nome

FROM estado e,

FROM fornecedor f

WHERE e.uf = f.uf;

\end{verbatim}


\choice \begin{verbatim}SELECT e.uf

FROM estado e

WHERE e.nome NOT IN (SELECT f.uf

                     FROM fornecedor f);

\end{verbatim}
\choice \begin{verbatim}SELECT e.nome

FROM estado e

WHERE e.uf IN (SELECT f.uf

               FROM fornecedor f);

\end{verbatim}


\choice \begin{verbatim}SELECT e.nome

FROM estado e

WHERE e.uf NOT IN (SELECT f.uf

                   FROM fornecedor f);

\end{verbatim}


\end{choices}

\question

Em uma lista circular duplamente encadeada com \( n \) elementos, o espaço ocupado apenas pelos apontadores é (assuma que um apontador ocupa \( p \) bytes):
\begin{choices}
\choice \( 4np \)
\choice \( 6np \)
\choice \( (np)^2 \)
\choice \( np \)
\choice \( 2np \)
\end{choices}

\question

Marque a alternativa onde todos os conceitos estão corretos.


\begin{choices}
\choice Nenhuma das alternativas anteriores.
\choice Em um diagrama de fluxo de dados uma entidade externa representa uma fonte ou destino das informações processadas pelo sistema e está fora dos limites do sistema modelado; cada processo pode ser refinado, para explicitar um maior detalhamento; um DFD pode conter vários níveis de detalhamento; um processo é um transformador de informação; o nível 0 de um DFD representa o sistema como um todo e indica as principais fontes e destinos das informações, usualmente referenciado por Diagrama de Contexto.
\choice Em um diagrama de fluxo de dados uma entidade externa representa um produtor ou um consumidor de informação e está fora dos limites do sistema modelado; cada processo deve ser refinado, para explicitar um maior detalhamento; um DFD pode conter vários níveis de detalhamento; um processo é um transformador de informação e também está fora do sistema; o nível 0 de um DFD representa o sistema como um todo e indica os principais usuários e as funções do sistema.
\choice Em um diagrama de fluxo de dados uma entidade externa representa uma fonte ou destino das informações processadas pelo sistema e está fora dos limites do sistema modelado; cada processo pode ser refinado, para explicitar um maior detalhamento; um DFD pode conter vários níveis de detalhamento; um processo é um transformador de informação e também está fora do sistema; o nível 0 de um DFD representa o sistema como um todo e indica as principais fontes e destinos das informações.
\choice Em um diagrama de fluxo de dados, uma entidade externa representa um produtor ou um consumidor de informação e está fora dos limites do sistema modelado; cada processo pode ser refinado, para explicitar um maior detalhamento; um DFD contém dois níveis de detalhamento; um processo é um transformador de informação e também está fora do sistema; o nível 0 de um DFD representa o sistema como um todo e indica os principais usuários e as funções do sistema.
\end{choices}

\question

Sobre a técnica conhecida como Z-buffer é correto afirmar que:
\begin{choices}
\choice É possível realizar o cômputo das variáveis envolvidas de forma incremental.
\choice As dimensões do Z-buffer são independentes das dimensões do frame buffer. 
\choice Nenhuma das alternativas acima está correta. 
\choice As primitivas geométricas precisam estar ordenadas de acordo com a distância em relação ao observador. 
\choice É uma técnica muito comum de detecção de colisão. 
\end{choices}

\question

Considere as seguintes tabelas em uma base de dados relacional:\\

Departamento (\underline{CodDepto}, NomeDepto)\\

Empregado (\underline{CodEmp}, NomeEmp, CodDepto, Salario)\\

Considere a seguinte consulta SQL:

\begin{verbatim}

SELECT d.CodDepto, d.NomeDepto, SUM(e.Salario)

FROM Departamento d

JOIN Empregado e ON d.CodDepto = e.CodDepto

GROUP BY d.CodDepto, d.NomeDepto

HAVING COUNT(e.CodEmp) > 2 AND AVG(e.Salario) > 40;

\end{verbatim}



Qual o resultado da consulta?
\begin{choices}
\choice Para cada departamento que tem mais que dois empregados e cuja média salarial, considerando todos empregados do departamento, exceto os dois primeiros, é maior que 40, obter o código de departamento, seguido do nome do departamento, seguido da soma dos salários dos empregados do departamento.
\choice Para cada empregado que tem mais que dois departamentos, ambos com média salarial maior que 40, obter o código de departamento, seguido do nome do departamento, seguido da soma dos salários dos empregados do departamento.
\choice Para cada departamento que tem mais que dois empregados e cuja média salarial é maior que 40, obter um grupo de linhas que contém, para cada empregado do departamento, o código de seu departamento, seguido do nome de seu departamento, seguido da soma dos salários dos empregados do departamento.
\choice Para cada departamento que tem mais que dois empregados e cuja média salarial é maior que 40, obter o código de departamento, seguido do nome do departamento, seguido da soma dos salários dos empregados do departamento.
\choice A consulta não retorna nada pois está incorreta.
\end{choices}

\question

O menor número possível de arestas em um grafo conexo com \( n \) vértices é:
\begin{choices}
\choice \( n \)
\choice \( n - 1 \)
\choice 1
\choice \( n^2 \)
\choice \( n/2 \)
\end{choices}

\question

Sejam os seguintes predicados de uma linguagem de primeira ordem:



\begin{itemize}

    \item \( N(x) \): \( x \) é número;

    \item \( P(x) \): \( x \) tem propriedade \( P \);

    \item \( x < y \): \( x \) é menor que \( y \).

\end{itemize}



E sejam os símbolos:

\begin{itemize}

    \item \( \forall \): quantificador universal;

    \item \( \Rightarrow \): operador se-então;

    \item \( \neg \): operador de negação.

\end{itemize}



Para a fórmula: \( \forall x (N(x) \Rightarrow \neg \forall y (N(y) \Rightarrow y < x)) \), qual alternativa abaixo \textbf{NÃO} constitui uma tradução possível?


\begin{choices}
\choice Para todo número, existe um outro número que é maior do que ele.
\choice Para qualquer \( x \), se \( x \) é número, então não é verdade que todos os números são menores do que ele.
\choice Não há um número tal que todos os números são menores do que ele.
\choice Não há um número menor do que outro número.
\choice Para todo número, não é verdade que qualquer número seja menor do que ele.
\end{choices}

\question

Ao analisar um novo sistema de gerenciamento de bancos de dados que surgiu no

mercado, você percebe que a entrada às operações do modelo de dados são tabelas

mas as saídas nunca são tabelas, são sempre outras estruturas. O vendedor desse

SGBD afirma que ele é um SGBD relacional. Pode-se afirmar que:


\begin{choices}
\choice Esse SGBD não é relacional pois no modelo relacional a única estrutura percebida pelos usuários são tabelas.
\choice Esse SGBD não é relacional pois, para ser verdadeiramente relacional, as operações do modelo de dados não podem trabalhar recebendo tabelas como entrada.
\choice Esse SGBD é relacional pois o vendedor afirmou que é.
\choice Esse SGBD é relacional pois consegue trabalhar com tabelas para a entrada das operações do modelo de dados. A saída não é importante.
\choice Nenhuma das alternativas anteriores.
\end{choices}

\question

(Adaptado de ENADE 2011) Considere o diagrama abaixo, que ilustra um banco de

dados que descreve os empregados que trabalham nos departamentos de uma empresa:

\begin{figure}[H]

  \centering

  \caption{Empregados e departamentos}

  \vspace{-0.2cm}

  \label{fig:empregados}

  %\fbox{

    \includegraphics[width=0.5\linewidth]{static/imgs/defaul.png}

  %}\\

  %\footnotesize{Fonte: xxxx.}

\end{figure}

\vspace{-0.4cm}

Na linguagem SQL, o comando mais simples para recuperar os códigos dos

departamentos que contém empregados que ganham menos do que 2000 e mais do que

5000 reais é:


\begin{choices}
\choice \begin{verbatim}SELECT cod_departamento

FROM empregados

WHERE salario <= 2000 AND salario >= 5000;

\end{verbatim}


\choice \begin{verbatim}SELECT cod_departamento

FROM empregados

WHERE salario BETWEEN 2000 AND 5000;

\end{verbatim}


\choice \begin{verbatim}SELECT cod_departametno

FROM empregados

WHERE salario <= 2000 OR salario >= 5000;

\end{verbatim}


\choice \begin{verbatim}SELECT cod_departamento

FROM empregados

WHERE salario < 2000 OR salario > 5000;

\end{verbatim}


\choice \begin{verbatim}SELECT cod_departamento

FROM empregados

WHERE salario < 2000 AND salario > 5000;

\end{verbatim}


\end{choices}

\question

Dentre as características do modelo relacional e do modelo de objetos em bancos de dados, qual afirmação é \textbf{INCORRETA}?


\begin{choices}
\choice O modelo de objetos é mais adequado para a representação de tipos abstratos de dados.
\choice O modelo de objetos provê uma representação bem próxima de linguagens de programação.
\choice O relacionamento binário \( N \times M \) é representado de modo semelhante nos dois modelos.	
\choice O relacionamento de herança é diretamente representado no modelo relacional.
\choice O modelo de objetos possui mais recursos estruturais para a representação de dados que o relacional.
\end{choices}

\question

Você estava ouvindo a discussão de dois colegas, o João e a Maria. Enquanto o

João afirmava que o conceito de banco de dados (BD) é extremamente diferente do

conceito de sistema de gerenciamento de bancos de dados (SGBD), a Maria

discordava e afirmava que o banco de dados nada mais é do que um componente do

SGBD. Nesse momento o professor entrou na sala, percebeu a discussão e falou

que:
\begin{choices}
\choice Não é possível saber quem estava certo ou errado pois não definiram qual o modelo de dados que estava sendo considerado na discusão.
\choice O João estava certo já que o banco de dados é, realmente, um conceito diferente do conceito do sistema de gerenciamento de bancos de dados, e existe de forma independente.
\choice Os dois estavam certos pois o conceito, na prática, depende da situação específica que se está considerando.
\choice Os dois estavam errados pois não são conceitos separados nem componentes, são conceitos sinônimos.
\choice A Maria estava certa já que o banco de dados é, realmente, um componente interno do sistema de gerenciamento de bancos de dados e não existe de forma independente.
\end{choices}

\question

A arquitetura ilustrada na figura abaixo é um exemplo típico de qual sistema?

(CU = unidade de controle; IS = \ingles{stream} de instruções; PU = unidade de

processamento; DS = \ingles{stream} de dados; LM = memória local)

\begin{figure}[H]

\begin{center}

\includegraphics[width=0.5\linewidth]{static/imgs/defaul.png}

\end{center}

\end{figure}

\vspace{-1cm}
\begin{choices}
\choice Clusters
\choice Acesso não uniforme à memória
\choice Multiprocessadores simétricos
\choice Processadores multicore
\choice Processadores seqüenciais
\end{choices}

\question

Uma integração de Sistemas Computacionais formando uma rede, tipicamente é implementada através da instalação de uma Arquitetura de Rede, que é composta de camadas e protocolos, em cada um dos elementos que compõem esta rede. Considere que estações “conversam” quando aplicações de usuários conseguem comunicar-se, sintática e semanticamente, através da Rede de Computadores. Baseados nesta premissa e em todos os conceitos associados à implementação e utilização das redes de computadores podemos afirmar como certo:


\begin{choices}
\choice Nenhuma delas é uma afirmação correta.
\choice Computadores com arquiteturas diferentes podem “conversar” através de multiplexadores.
\choice Computadores com arquiteturas de rede parecidas conseguem “conversar”.
\choice Computadores com arquiteturas de redes diferentes podem “conversar” através de um \textit{gateway} ou conversor de protocolos.
\choice Computadores com arquiteturas de redes diferentes conseguem “conversar”.
\end{choices}

\question

(ENADE 2017) No modelo relacional de dados, uma junção entre tabelas cria uma

outra tabela derivada de duas ou mais tabelas de acordo com os critérios

especificados para a junção, de modo similar às regras da teoria dos

conjuntos. Nos conjuntos a seguir, as áreas hachuradas representam os resultados

das consultas SQL em que ``id' é uma chave primária e ``fk' é uma chave

estrangeira:

\begin{figure}[H]

  \centering

  \includegraphics[width=0.5\linewidth]{static/imgs/defaul.png}

\end{figure}

Assinale a opção em que a declaração SQL é equivalente ao conjunto apresentado a

seguir:

\begin{figure}[H]

  \centering

  \includegraphics[width=0.5\linewidth]{static/imgs/defaul.png}

\end{figure}


\begin{choices}
\choice \begin{verbatim}SELECT *

FROM tabela1 t1

INNER JOIN tabela2 t2 ON (t1.id = t2.fk);

\end{verbatim}


\choice \begin{verbatim}SELECT *

FROM tabela1 t1 WHERE EXISTS (SELECT 1

                              FROM tabela2 t2

                              WHERE t2.id = t1.fk);

\end{verbatim}


\choice \begin{verbatim}SELECT *

FROM tabela1 t1

RIGHT JOIN tabela2 t2 ON (t1.id = t2.fk)

WHERE t1.id IS NULL;

\end{verbatim}


\choice \begin{verbatim}SELECT *

FROM tabela1 t1

RIGHT JOIN tabela2 t2 ON (t1.id = t2.fk)

WHERE t2.fk <> NULL;

\end{verbatim}


\choice \begin{verbatim}SELECT *

FROM tabela1 t1

WHERE NOT EXISTS (SELECT 1

                  FROM tabela1 t2

                  WHERE t2.id = tk.fk);

\end{verbatim}


\end{choices}

\question

Avalie o banco de dados de professores e alunos, mostrado a seguir. Esse banco

de dados pretende armazenar todas as turmas, sendo que um professor pode dar

aula para vários alunos ao mesmo tempo, e um mesmo aluno pode ter aula com

diferentes professores ao mesmo tempo.

\begin{figure}[H]

  \centering

  \caption{Professores, turmas e alunos}

  \label{fig:professores}

  %\fbox{

    \includegraphics[width=0.5\linewidth]{static/imgs/defaul.png}

  %}\\

  %\footnotesize{Fonte: xxxx.}

\end{figure}

Esse banco de dados resolveráo problema de armazenar os dados das turmas?
\begin{choices}
\choice Não é possível determinar somente avaliando o diagrama exibido, seriam necessárias mais informações do dicionário de dados.
\choice Não, pois as cardinalidades e os sentidos dos relacionamentos entre 'professores' e 'turmas', e entre 'turmas' e 'alunos' não cria um relacionamento N:N entre 'professores' e 'alunos'.
\choice Não, pois o relacionamento entre 'turmas' e 'professores' está errado: a PK 'cod\_turma' deveria ter ido para a tabela 'professores' como uma FK, mas no diagrama está ao contrário.
\choice Sim, pois o diagrama mostra claramente que um professor pode ter zero ou mais turmas, e cada aluno pode participar de zero ou mais turmas também.
\choice Sim, pois as cardinalidades e os sentidos dos relacionamentos entre 'professores' e 'alunos' indica, no final, um relacionamento N:N entre essas duas tabelas.
\end{choices}

\question

(Adaptado do ENADE 2005) Analise o relacionamento existente entre as tabelas

``clubes' e ``jogadores', nas alternativas abaixo, e marque a alternativa que

indica que todo jogador deve pertencer a um único clube, e todo clube poder ter

ou não jogadores.
\begin{choices}
\choice \includegraphics[width=0.5\linewidth]{static/imgs/defaul.png} \vspace{0.5cm}
\choice \includegraphics[width=0.5\linewidth]{static/imgs/defaul.png} \vspace{0.5cm}
\choice \includegraphics[width=0.5\linewidth]{static/imgs/defaul.png} \vspace{0.5cm}
\choice \includegraphics[width=0.5\linewidth]{static/imgs/defaul.png} \vspace{0.5cm}
\choice \includegraphics[width=0.5\linewidth]{static/imgs/defaul.png} \vspace{0.5cm}
\end{choices}

\question

Considere as seguintes afirmações sobre redes neurais artificiais: \begin{enumerate}[label=\Roman*.] \item Um perceptron elementar só computa funções linearmente separáveis. \item Não aceitam valores numéricos como entrada. \item O "conhecimento" é representado principalmente através do peso das conexões. \end{enumerate} São corretas:
\begin{choices}
\choice Apenas II e III
\choice Apenas I e III
\choice Apenas III
\choice Apenas I e II 
\choice I, II e III
\end{choices}

\question

Qual das seguintes igualdades é verdadeira?



\begin{enumerate}

    \item[(I)] \( n^2 = \mathcal{O}(n^3) \)

    \item[(II)] \( 2 \cdot n + 1 = \mathcal{O}(n^2) \)

    \item[(III)] \( n^3 = \mathcal{O}(n^2) \)

    \item[(IV)] \( 3 \cdot n + 5 \cdot n \log n = \mathcal{O}(n) \)

    \item[(V)] \( \log n + \sqrt{n} = \mathcal{O}(n) \)

\end{enumerate}



As opções para a resposta correta são:
\begin{choices}
\choice somente I, III e IV
\choice somente II, III e IV
\choice somente I e II
\choice somente III, IV e V
\choice somente I, II e V
\end{choices}

\question

Em SQL, qual comando é utilizado para combinar linhas de duas ou mais tabelas?


\begin{choices}
\choice COMBINE
\choice MERGE
\choice LINK
\choice CONNECT
\choice JOIN
\end{choices}

\question

São linguagens clássicas que representam aplicações científicas, aplicações 

comerciais, aplicações de inteligência artificial, e aplicações gerais,

respectivamente:
\begin{choices}
\choice Fortran, COBOL, LISP, C
\choice MATLAB, COBOL, SML, LISP
\choice COBOL, Java, LISP, Fortran
\choice R, COBOL, C, Python
\choice Java, MATLAB, Scheme, JavaScript
\end{choices}

\question

Em um banco de dados relacional, o que é uma chave primária?


\begin{choices}
\choice É denominada por FK
\choice É qualquer coluna ou conjunto de colunas que podem servir como índice
\choice É uma validação que evita o cadastro de dados errados
\choice É um atributo (ou um conjunto de atributos) que permite identificar única e exclusivamente cada registro da tabela
\choice Nenhuma das respostas acima
\end{choices}

\question

Qual é o número mínimo de comparações necessário para encontrar o menor elemento de um conjunto qualquer não ordenado de \( n \) elementos?
\begin{choices}
\choice \( n + 1 \)
\choice \( n \log n \)
\choice 1
\choice \( n \)
\choice \( n - 1 \)
\end{choices}

\question

Em relação ao paradigma de programação cliente-servidor. Qual das afirmativas abaixo é FALSA?
\begin{choices}
\choice Um aplicativo cliente é um programa arbitrário que se torna temporariamente um cliente quando for necessário o acesso remoto a um serviço, mas também executa processamento local.
\choice Um aplicativo servidor é um programa de propósito especial dedicado a fornecer um serviço, mas pode tratar de múltiplos clientes remotos ao mesmo tempo. 
\choice Um aplicativo cliente pode acessar múltiplos serviços quando necessário.
\choice Um aplicativo servidor aceita contato de clientes arbitrários, mas oferece um único serviço. 
\choice Um aplicativo servidor inicia ativamente o contato com clientes 
\end{choices}

\question

Por que não podemos utilizar apelidos de coluna na cláusula WHERE da SQL?
\begin{choices}
\choice Porque a cláusula WHERE é avaliada antes da SELECT.
\choice Porque os apelidos só são avaliados depois da cláusula ORDER BY.
\choice Nós podemos utilizar apelidos na cláusula WHERE.
\choice Porque a cláusula SELECT é avaliada antes da WHERE.
\choice Porque só podemos utilizar aliases, e não apelidos.
\end{choices}

\question

Seja a seguinte linguagem, onde \(\epsilon\) representa o string vazio e \$ representa um marcador de fim de entrada:



$S \rightarrow ABCD \\$\\

$A \rightarrow a \mid \epsilon \\$\\

$B \rightarrow a \mid \epsilon \\$\\

$C \rightarrow c \mid \epsilon \\$\\

$D \rightarrow S \mid c \mid \epsilon$\\



É incorreto afirmar que:
\begin{choices}
\choice O conjunto FOLLOW(D) é igual a FOLLOW(S)
\choice O conjunto FOLLOW(B) = c, \$
\choice O conjunto FIRST(D) é igual ao conjunto FIRST(S)
\choice O conjunto FIRST(A) = a, \(\epsilon\)
\choice O conjunto FOLLOW(A) = a, c, \$
\end{choices}



    %%%%%%%%%%%%%%%%%%%%%%%%%%%%%%%%%%%%%%%%%%%%%%%%%%%%%%%%%%%%%%%%%%%%%%%%%%%%%%%%
    %%% Finaliza o bloco de questões:
    \end{questions}
    \end{document}
    